\documentclass[12pt]{article}

\usepackage[a4paper, total={6in, 8in}]{geometry}
\usepackage{textcomp}

\begin{document}
\setlength{\parindent}{0pt}

\def \t {\textrightarrow}
\def \v {\vspace{1em}}
\def \bi {\begin{itemize}}
\def \ei {\end{itemize}}
\def \s[#1] {\section*{#1}}
\def \ss[#1] {\subsection*{#1}}
\def \sss[#1] {\subsubsection*{#1}}

\s[Il vento scrive - D'Annunzio]
Catullo scrive un carme "Parole scritte nell'acqua e nel vento" \t sono le parole d'amore, che vengono cancellate subito e vengono portate via dal vento \\
Ama quindi la cultura classica \\
Sempre parte di Alcyone, la stagione estiva \t insieme agli altri due che sono "La sabbia del tempo" e "Nella belletta" \\
Sono madrigali d'estate \t metrica del madrigale vicina a quella del sonetto \t due terzine e poi strofa finale con due versi \\
Tema della natura-uomo, uomo-natura \\
Immilli = dantismo \t nel senso di comparire \\
Il sorriso viene cancellato dalla docile sabbia (sinestesia) \t e il vento scrive con le penne delle ali dei volatili \\
Favella è di solito usato come verbo \t le parole del vento in realta parlerebbero \\
Parlare i segni \t usato in modo transitivo, che è insolito \\
Quando il sole declina \t crea la luce che si riflette creando la spaccatura con i raggi \t immagine che si trova anche nella sabbia del tempo, con l'ombra della tacita meridiana \\
Come se fosse l'ombra delle ciglie sulle gote \t di nuovo figura femminile, e sensaulismo \\
È presente coincidenza metrico-sintattica \t le strofe sono chiuse \\
S'immilli è dantismo, ma anche neologismo \t e anche iperbole, vuol dire moltiplicarsi mille volte

\v

Il vento che scrive, la vanita delle cose e del tutto \t l'aridità della vita, e gli oggetti incarnano significati profondi che riportano a punti fermi dell'autore: il tempo che passa, e nello specifico il paradosso \\
L'immagine di lei da un lato sparisce per il vento, ma rimane immillato \t paradosso, infinito e zero \\
Spirito decadente, che conferma la passione di d'annunzio per la natura che è essere umano \t chiaro ripiegamento intimistico

\ss[Nella belletta]
Fango lutulento \t iperbato in cui cuoce sta alla fine \\
Dolcigna afa di morte \t odore dolciastro nauseabondo \\
La rana con le bolle d'aria che salgono in silenzio
\end{document}
